\section{Local stretching with a signed drift}
\label{sec:stretching}

We now enlarge the distance between two closed sets while retaining the
quadratic reserve constructed in Proposition~\ref{pre:reserve}.  The enlargement
must have compact support: this will allow infinitely many applications on
a noncompact manifold.  We adapt the local graph construction of
\cite[Theorem~3.1]{OpenAIWidth}, using Proposition~\ref{pressure:law} to allow the drift
to take either sign.

The loss in the following statement has two parts.  One is charged to a
decrease of a metric trace; these trace decreases telescope in the
iteration.  The other can be chosen arbitrarily small.  We write $a_+=\max\{a,0\}$ for the positive
part of a real number.

\begin{lemma}[Local stretching]\label{stretch:local}
Let $n\ge4$ and $0<\delta<1/10$.  There are constants $W\ge1$ and $T\ge10$,
depending only on $n,\delta$, with the following property.  Let $M$ be a
connected smooth $n$-manifold without boundary, let $h_0$ be a complete
smooth metric on $M$, let $h\ge h_0$ be a smooth metric, and let
$t\in C^\infty(M,\R)$.  Suppose that $A\subset M$ is nonempty and compact,
that $B\subset M$ is nonempty and closed, and that
$\dist_{h_0}(A,B)\ge W$.  Put
\[
 V(t)=t^2+T^2,\qquad L_*=1.004.
\]
For every $0<d<\delta$, there exist a smooth complete metric $h'\ge h$ and
a smooth function $t'\ge t$, whose changes from $h,t$ have compact support
in $\{x:\dist_{h_0}(A,x)<W\}$, such that
\begin{align}
 \Dcal(h',t')\ge{}&\Dcal(h,t)+(V(t')-V(t))_+
 \notag\\[-2pt]
 &-\delta\bigl(\tr_hh_0-\tr_{h'}h_0\bigr)V(t')-dV(t'),
 \label{stretch:conclusion}\\
 \dist_{h'}(A,B)\ge{}&L_*W.
 \label{stretch:distance-conclusion}
\end{align}
\end{lemma}

The proof constructs a graph that follows a large height function near
$A$ and becomes small before leaving its $W$-neighborhood.  At appreciable
slope, the graph metric increases lengths.  At small slope, a modified
graph metric makes the curvature error arbitrarily small.  After solving
the graph equation on a compact extension, we estimate its scalar
curvature and then cut off the change.

\subsection{Slope modification and spatial cutoffs}
\label{stretch:setup}

Fix the constants
\begin{equation}\label{stretch:constants}
 \begin{gathered}
 s_0=\frac1{100},\qquad A_0=1000,\qquad P_0=20,
 \qquad C_* =s_0^{-1},\\
 b_0=\frac{\delta s_0^2}{2000},\qquad
 \epsilon=\min\left\{\frac{s_0}{100},\frac1{8n},
                              \frac1{8C_*}\right\}.
 \end{gathered}
\end{equation}
Apply Proposition~\ref{pressure:law} with these parameters, obtaining
$S_0,T,F$.  Throughout this section,
$N(t)=\sqrt{t^2+T^2}$ and $V(t)=N(t)^2$.
We may take $W$ to be a numerical constant at least $10^5$; this size
allows the cutoffs below and the final length comparison.

Let $0<a\le1/2$.  This parameter will be fixed small enough in terms of
the prescribed loss $d$.  Choose a smooth nondecreasing function
$\sigma:[0,\infty)\to[0,1]$ that is zero on $[0,s_0]$, one on
$[5s_0,\infty)$, and satisfies $0\le\sigma'\le60$.  For $s\ge0$ set
\begin{equation}\label{stretch:slope-functions}
 \mu=a^2+(1-a^2)\sigma(s),\qquad
 \theta=\frac{\mu s^2}{1+s^2},\qquad
 l=\sqrt{1-\theta},\quad \beta=\log l,
 \quad k_1=\sqrt\theta.
\end{equation}
Thus the graph modification of Lemma~\ref{pre:graph-comparison} is the
ordinary graph metric when $s\ge5s_0$, whereas its size is reduced by
$a^2$ near zero slope.  For the graph equation, we use the multiplier
\[
 j(s)=\sqrt{\theta+a^2s_0^2}.
\]
It is smooth as a function of the gradient covector, and $j(0)=as_0>0$.
This keeps the derivative of the pressure with respect to height positive
at critical points of the graph.  The same multiplier will also make the
small-slope curvature error proportional to $a^2$.

We next arrange the spatial support.  Write
$r_A(x)=\dist_{h_0}(A,x)$.  Choose smooth compactly supported functions
$\psi\ge0$ and $\eta_1,\zeta\in[0,1]$ such that
\begin{equation}\label{stretch:cutoffs}
 \begin{gathered}
 \inf_A\psi>W/5,\qquad |\dd\psi|_{h_0}\le1,
       \qquad \supp\psi\subset\{r_A<W/2\},\\
 \eta_1=1\text{ on }\{r_A\le.55W\},\qquad
 \eta_1=0\text{ on }\{r_A\ge.8W\},
       \qquad |\dd\eta_1|_{h_0}\le s_0/8,\\
 \zeta=1\text{ on }\{r_A\le.82W\},\qquad
 \zeta=0\text{ on }\{r_A\ge.98W\}.
 \end{gathered}
\end{equation}
These cutoffs require only completeness of $h_0$.  Indeed, finite-radius
closed neighborhoods of the compact set $A$ are compact.  Start with
Lipschitz functions of $r_A$, leaving positive buffers around each
specified constant region.  In particular, smooth
$(W/3-r_A)_+$ with value error less than $W/100$ and gradient bound
$1.05$, then divide by $1.05$ to obtain $\psi$.  The profile for $\eta_1$
can be smoothed with gradient at most $12/W\le s_0/8$.
To carry out the smoothing, use convolutions in sufficiently small charts
on the compact neighborhood, where metric distortion is as small as
desired, and a partition of unity.  Choose the local approximation errors
small enough to control the partition derivatives.  Nonnegative kernels
preserve the range, and the buffers preserve the stated constant regions.
The same procedure supplies $\zeta$.

For $0<\rho\le1/2$, put
\begin{equation}\label{stretch:eta}
 \eta=\eta_\rho=\rho+(1-\rho)\eta_1^2.
\end{equation}
Thus $\eta=1$ near $\supp\psi$, while $\eta=\rho$ outside
$\{r_A<.8W\}$.  Moreover, since $h\ge h_0$,
\begin{equation}\label{stretch:eta-bounds}
 |\dd\eta|_h\le s_0/4,\qquad
 \frac{|\dd\eta|_h^2}{\eta}\le1.
\end{equation}
The square in \eqref{stretch:eta} gives the second estimate uniformly
as $\rho\downarrow0$.  For each $\rho$ we will construct a smooth graph
height $u_\rho$ close to $\psi$.  The arrangement in Figure~\ref{stretch:figure}
leaves a region where the graph tends to zero before the final cutoff
$\zeta$ begins to vary.

\begin{figure}[tbp]
\centering
\begin{tikzpicture}[x=.98cm,y=.8cm,font=\small]
 \fill[gray!9] (7,-.4) rectangle (11,2.55);
 \foreach \x in {4,5,7,8,10,11}
   \draw[gray!55,densely dotted] (\x,-.55)--(\x,2.5);
 \fill[orange!22] (0,1.8) rectangle (3.8,2.35);
 \draw (0,1.8) rectangle (3.8,2.35);
 \node at (1.9,2.075) {$\supp\psi$};
 \node[anchor=east] at (-.25,2.075) {target support};
 \fill[blue!15] (0,.8) rectangle (5,1.35);
 \fill[pattern=north east lines,pattern color=blue!45]
                  (5,.8) rectangle (7,1.35);
 \draw (0,.8) rectangle (7,1.35);
 \draw (5,.8)--(5,1.35);
 \node at (2.5,1.075) {$\eta_1=1$};
 \node at (9,1.075) {$\eta_1=0,\quad\eta=\rho$};
 \fill[blue!15] (0,-.2) rectangle (8,.35);
 \fill[pattern=north east lines,pattern color=blue!45]
                  (8,-.2) rectangle (10,.35);
 \draw (0,-.2) rectangle (10,.35);
 \draw (8,-.2)--(8,.35);
 \node at (4,.075) {$\zeta=1$};
 \node[anchor=east] at (-.25,.075) {final cutoff};
 \draw[-{Stealth[length=2mm]}] (0,-.65)--(11.6,-.65);
 \foreach \x/\lab in {0/0,4/.50,5/.55,7/.80,8/.82,10/.98,11/1}
   {\draw (\x,-.60)--(\x,-.72);
    \node[below] at (\x,-.72) {$\lab$};}
 \node[anchor=west] at (11.65,-.65) {$r_A/W$};
 \draw[<->] (7,2.8)--(11,2.8);
 \node[above] at (9,2.8) {$u_\rho\longrightarrow0$ locally};
\end{tikzpicture}
\caption{The spatial cutoffs, with $\eta=\rho+(1-\rho)\eta_1^2$.
The first bar shows the containment $\supp\psi\subset\{r_A<W/2\}$;
hatching indicates where a cutoff may vary.  The graph heights $u_\rho$
constructed below tend to zero in $C^3$ locally outside
$\{r_A\le.8W\}$ as $\rho\downarrow0$, leaving room for the final metric
and drift cutoff.  Spacing is schematic; endpoints are exact multiples of $W$.}
\label{stretch:figure}
\end{figure}

\subsection{Solving the graph equation and obtaining locality}
\label{stretch:equation}

Work for the moment on a closed auxiliary manifold extending a
neighborhood of $\{r_A\le1.1W\}$.  If $M$ is compact, use $M$ itself.
Otherwise, enclose this compact neighborhood in a relatively compact
smooth domain, modify the data to product form near its boundary, and
double.  Extend $h,t$ smoothly and the compactly supported cutoffs by
zero.  Define $\eta$ from the extended $\eta_1$ as above.  This auxiliary
manifold is fixed as $\rho$ tends to zero; all geometric constants in this
subsection may depend on it.

For $0\le\lambda\le1$, consider the equation
\begin{equation}\label{stretch:pde}
 \begin{gathered}
 \diver_{h,t}\mathcal A(\dd u)
     =j(s)F\left(\frac{u-\lambda\psi}{\eta},t\right),\\
 \mathcal A(\dd u)=\frac{\sqrt\mu}{\sqrt{1+s^2}}\nabla_hu,
 \qquad j(s)=\sqrt{\theta+a^2s_0^2},
 \qquad s=|\dd u|_h.
 \end{gathered}
\end{equation}
We solve it within the open domain
$|u-\lambda\psi|/\eta<S_0$.  Both $\mathcal A$ and $j$ are smooth
functions of the covector $\dd u$, including at zero.  The flux derivative
has tangential and radial eigenvalues $k_1/s$ and $(k_1)_s$, respectively;
both are positive for $s>0$ and tend to $a$ at zero.

We first obtain estimates for every smooth solution in this domain.
There is $S_1<S_0$, independent of $\rho,\lambda$, such that
\begin{equation}\label{stretch:height}
 |u-\lambda\psi|\le S_1\eta.
\end{equation}
To see this, compare with $\lambda\psi\pm S\eta$ for $S<S_0$.
Their first two derivatives are uniformly bounded as $S\uparrow S_0$,
whereas $j\ge as_0$ and the poles of $F$ are uniform on the compact range
of $t$.  For $S$ sufficiently close to $S_0$, these functions are strict
upper and lower barriers.  At a positive maximum of a violation, the
gradients agree, ellipticity orders the second derivatives, and $F_x>0$
orders the right sides, a contradiction.  The same fixed $S$ therefore
works for all $\rho,\lambda$.

There is also a gradient bound independent of $\rho,\lambda$.  At a point
of maximum slope, there is nothing to prove if $s<1$.  Otherwise $\mu=1$
in a neighborhood, and \eqref{stretch:pde} is the ordinary graph equation
of Lemma~\ref{pre:angle}, with
\[
 p(x,z,\omega)=\sqrt{1-\omega^2+a^2s_0^2}\,
 F\left(\frac{z-\lambda\psi(x)}{\eta(x)},t(x)\right),
 \qquad \omega=(1+s^2)^{-1/2}.
\]
By \eqref{stretch:height}, the values of $F,F_t,F_x$ range in a compact
set uniformly in $\rho,\lambda$.  Holding $\omega$ fixed and writing
$\xi=(u-\lambda\psi)/\eta$, we have on the graph
\[
 \partial_\nu\left(\frac{z-\lambda\psi}{\eta}\right)
 =\frac{\nu_z-\lambda\,\dd\psi(\nu_h)
                         -\xi\,\dd\eta(\nu_h)}{\eta}.
\]
Thus the zeroth order coefficient of the angle operator is bounded below
by $-C/\eta$, after discarding the nonnegative term $|\mathrm{II}|^2$.
Here $C$ is independent of $\rho,\lambda$.  On the other hand,
\[
 p_z=\frac{jF_x}{\eta}\ge\frac{A_0T}{\sqrt2\,\eta}
\]
at the point under consideration.  At this minimum of $\omega$, the
angle equation gives
\[
 -p_z=\mathcal L\omega\ge-\frac{C}{\eta}\omega.
\]
It follows that $\omega$ is uniformly bounded below, as required.

For each fixed $\rho>0$, these estimates confine the solutions to compact
jet ranges strictly inside the domain of the equation.  The elliptic
estimates described in Section~\ref{pre:elliptic-tools} give all higher
bounds.  The linearization of the left side minus the right side has
zeroth order coefficient $-jF_x/\eta<0$, so it is invertible on the closed
manifold by the maximum principle and linear elliptic theory.  The method
of continuity therefore applies from the solution $u=0$ at $\lambda=0$
to $\lambda=1$.  Write $u_\rho$ for the resulting smooth solution.
The higher estimates used for this continuation may depend on $\rho$;
the height and gradient estimates do not.

The graph equation has now produced a height close to $\psi$.  Its next
property is what permits a compactly supported change of metric:
\begin{equation}\label{stretch:locality}
 u_\rho\longrightarrow0\quad\text{in }C^3
 \text{ on compact subsets of any open region where }
                       \psi=\eta_1=0.
\end{equation}
In such a region $\eta=\rho$.  Expanding \eqref{stretch:pde} in a
coordinate chart gives
\begin{equation}\label{stretch:reaction}
 a^{ij}\partial_i\partial_j u_\rho+b^i\partial_i u_\rho
                         -\frac{c_\rho}{\rho}u_\rho=0,
 \qquad
 c_\rho=j\int_0^1F_x(qu_\rho/\rho,t)\,\dd q.
\end{equation}
The previously established gradient bound makes $(a^{ij})$ uniformly
elliptic and bounded and makes $b^i$ bounded.  The first order terms have
this form because the flux vanishes at zero covector.  By
\eqref{stretch:height} and Proposition~\ref{pressure:law}, $c_\rho$ is
bounded above and below by positive constants, independently of $\rho$.

On a coordinate box with closure in this region, let $\ell_f$ be the
inward affine distance from its face $f$.  For sufficiently small fixed
$\gamma>0$, the function
\[
 C\sum_f\exp(-\gamma\ell_f/\sqrt\rho)
\]
is a supersolution of the operator in \eqref{stretch:reaction}, for all
small $\rho$.  Indeed, its second derivative terms are bounded by a
constant times $\gamma^2/\rho$, its first derivative terms by a constant
times $\gamma/\sqrt\rho$, and its negative zeroth order term dominates
both.  Choose $C$ to dominate the uniform height bound on the faces and
compare with $u_\rho$ and $-u_\rho$.  This proves exponential decay of
the height on every interior subbox.

To turn this into derivative decay, rescale around an interior point by
$x=x_0+\rho y$ and $u_\rho(x)=\rho v(y)$.  The equation takes the
divergence form
\[
 \partial_{y_i}\mathfrak A^i(x_0+\rho y,Dv)
           =\rho\,\mathfrak B(x_0+\rho y,v,Dv).
\]
Here $|v|\le S_1$ by \eqref{stretch:height}, while $Dv=Du_\rho$ is
uniformly bounded.  All coefficient derivatives on these ranges are
uniformly bounded, and ellipticity is uniform.  The interior estimates
from Section~\ref{pre:elliptic-tools} give bounded derivatives of $v$
of every order on smaller fixed balls.  Returning to the original
coordinates yields bounds $C_j\rho^{1-j}$ for derivatives of order $j$.
For fixed nested subboxes $K\subset\operatorname{int}K'$, interior
interpolation gives
\[
 \|u_\rho\|_{C^3(K)}
 \le C\|u_\rho\|_{C^0(K')}^{1/4}
          \|u_\rho\|_{C^4(K')}^{3/4}
       +C\|u_\rho\|_{C^0(K')}.
\]
The $C^4$ norm is at most $C\rho^{-3}$, while the $C^0$ norm decays
exponentially in $1/\sqrt\rho$.  The right side tends to zero, proving
\eqref{stretch:locality}.

\subsection{Curvature balance}
\label{stretch:curvature}

For the slope functions \eqref{stretch:slope-functions},
\[
 0\le-s\beta_s\le2.
\]
Indeed, $-s\beta_s=s\theta_s/[2(1-\theta)]$; on the switching interval
use $s\le.05$ and $\sigma'\le60$, and outside it the formula is immediate.
Thus Lemma~\ref{pre:graph-comparison} applies to $u=u_\rho$.  Denote its
output by
\begin{equation}\label{stretch:raw-data}
 h_u=h+\frac{\theta}{(1-\theta)s^2}\,\dd u^2,
                 \qquad t_u=t-\beta.
\end{equation}
Here $\dd u^2$ denotes $\dd u\otimes\dd u$.
These expressions are smooth across $s=0$: the coefficient of $\dd u^2$
extends as $\mu/[1+(1-\mu)s^2]$.  All comparisons with $h_0$ below take
place in the original neighborhood of $A$; no extension of $h_0$ is
needed.

On $\{s>0\}$, put $e=\nabla_hu/s$ and $b=\sqrt{h_0(e,e)}\le1$.
The rank-one inverse metric formula gives
\begin{equation}\label{stretch:trace}
 \tr_hh_0-\tr_{h_u}h_0=\theta b^2.
\end{equation}
Use $J,H,K,M_1$ as in Lemma~\ref{pre:graph-comparison}: $J$ is the
second fundamental form of the levels of $u$, $H=\tr J$,
$K=-e\log k_1$, and $M_1=H-et-K$.  In particular,
\[
 \Dcal(h_u,t_u)-\Dcal(h,t)\ge\theta E,\qquad
 E=|J|^2+M_1^2+K^2/l^2-2M_1K+2eM_1.
\]
Increasing the drift by $-\beta$ has a controlled effect on the reserve:
\begin{equation}\label{stretch:reserve-increment}
 \frac{(V(t-\beta)-V(t))_+}{\theta}
       \le4(1+s)N(t),\qquad s>0.
\end{equation}
Indeed, if $b_*=-\beta\ge0$, the numerator is at most
$2N(t)b_*+b_*^2$.  When $\theta\le1/2$ one has $b_*\le\theta$.
When $\theta>1/2$, one has $s>1$ and
$b_*\le\tfrac12\log(1+s^2)\le\sqrt s$; these inequalities give
\eqref{stretch:reserve-increment}.

Together with \eqref{stretch:trace}, this identifies the estimate needed
at slopes $s\ge s_0$: it suffices to prove
$E\ge4(1+s)N(t)-\delta b^2V(t_u)$.  Below $s_0$, we will instead use
the adjustable small error in Lemma~\ref{stretch:local}.

To express the graph equation in the normal variables, put, for $s>0$,
\begin{equation}\label{stretch:chi}
 \chi=\frac{j}{k_1}=\sqrt{1+\frac{a^2s_0^2}{\theta}},\qquad
 m_1=\frac{\dd\log\chi}{\dd\log k_1}
       =-\frac{a^2s_0^2}{\theta+a^2s_0^2}.
\end{equation}
The factor $\chi$ normalizes the pressure after division by $k_1$, and
$m_1$ records its derivative.  Directly from these formulas,
\begin{equation}\label{stretch:slope-bounds}
 \begin{gathered}
 -1\le m_1\le0,\qquad 1\le\chi\le2\quad(s\ge s_0),\\
 m_1\le-\tfrac12,\qquad \chi\ge s_0/s,\qquad
 \theta\chi^2\le2a^2s_0^2\quad(0<s\le s_0).
 \end{gathered}
\end{equation}
For $s\le s_0$, substitute $\theta=a^2s^2/(1+s^2)$; for $s\ge s_0$,
use $\theta\ge a^2s^2/(1+s^2)$ to bound $\chi$.

Equation~\eqref{stretch:pde} at $\lambda=1$ says
$M_1=\chi F$, with $F$ evaluated at $(\xi,t)$ and
$\xi=(u-\psi)/\eta$.  Since $e\log\chi=-m_1K$ and
$et=H-M_1-K$, differentiation gives
\[
 eM_1=-m_1KM_1+
       \chi\bigl(F_t(H-M_1-K)+F_xe\xi\bigr).
\]
Completing the $H$- and $K$-squares now yields
\begin{equation}\label{stretch:square-estimate}
 \begin{split}
 E&\ge\chi^2\mathcal C+2\chi F_xe\xi,\\
 \mathcal C&=
 [1-l^2(1+m_1)^2]F^2
 -2[1+l^2(1+m_1)]FF_t-nF_t^2.
 \end{split}
\end{equation}
For clarity, the discarded quantity is exactly
\[
 |J|^2-\frac{H^2}{n-1}
 +\frac{(H+(n-1)\chi F_t)^2}{n-1}
 +\left(\frac K l-l\chi((1+m_1)F+F_t)\right)^2
 +\theta\chi^2F_t^2,
\]
which is nonnegative.  The choice of $j$ now has two useful effects.
At small slope, $m_1\le-1/2$ leaves the coefficient of $F^2$ at least
$3/4$, while $\theta\chi^2=O(a^2)$ makes its eventual error small.
At larger slope, the two alternatives in the pressure law supply the
balance: a large value of $F_x$ is accompanied by a favorable term
$-FF_t$.

\paragraph{Slopes at least $s_0$.}
The pressure bounds imply $F_t^2\le\epsilon|FF_t|$.  Since
$n\epsilon\le1/8$ and $0\le1+m_1\le1$, they give
\begin{equation}\label{stretch:large-C}
 \mathcal C\ge-5\epsilon F_x,\qquad
 -FF_t\ge P_0F_x\ \Longrightarrow\ \mathcal C\ge P_0F_x.
\end{equation}
Indeed, if $FF_t\ge0$, its coefficient has absolute value at most $4$;
if $FF_t<0$, that coefficient contributes at least $2|FF_t|$, whereas
the last term costs at most $n\epsilon|FF_t|$.

If $\eta\ne1$, then $\psi=0$ in a neighborhood, so
\[
 e\xi=\frac{s-\xi e\eta}{\eta}\ge\frac{3s}{4\eta},
\]
using $|\xi|\le S_1<1$, \eqref{stretch:eta-bounds}, and $s\ge s_0$.
Equations~\eqref{stretch:square-estimate}--\eqref{stretch:large-C}
therefore give
$E\ge(3s/2-20\epsilon)F_x\ge sF_x$.
If $\eta=1$, then $\dd\eta=0$ and
$e\xi\ge s-b$ by $|\dd\psi|_{h_0}\le1$.
For $s>8b$ the same estimates give
$E\ge(7s/4-20\epsilon)F_x\ge sF_x$.

It remains to consider $s_0\le s\le8b$ with $\eta=1$.
Here $2\chi(s-b)-s\ge-4$, so $E\ge sF_x-10F_x$.
On the pressure alternative $-FF_t\ge P_0F_x$, the positive term
$\chi^2P_0F_x$ more than pays this loss and gives $E\ge sF_x$.
Otherwise $F_x\le b_0N(t)^2$, and the deficit is at most
$10b_0N(t)^2$.  Since $s\le8$, we have $0\le-\beta<3$, and $T\ge10$
implies
\[
 N(t)^2\le2V(t_u).
\]
This comparison is valid for signed $t$: for $0\le c<3$,
$2V(t+c)-V(t)=(t+2c)^2+T^2-2c^2>0$.
Moreover, $b\ge s_0/8$, and hence
\[
 10b_0N(t)^2\le20b_0V(t_u)
 =\frac{\delta s_0^2}{100}V(t_u)
 \le\delta b^2V(t_u).
\]
We have proved throughout $\{s\ge s_0\}$ that
\begin{equation}\label{stretch:large-slope}
 E\ge sF_x-\delta b^2V(t_u)
   \ge4(1+s)N(t)-\delta b^2V(t_u),
\end{equation}
where the last inequality uses $F_x\ge1000N(t)$ and $s\ge.01$.
Together with \eqref{stretch:reserve-increment} and
\eqref{stretch:trace}, this proves the required curvature inequality
without the small loss $d$ on this slope range.

\paragraph{Slopes below $s_0$.}
Here $1+m_1\le1/2$.  The pressure estimate $|F_t|\le\epsilon|F|$
therefore gives
\[
 \mathcal C\ge(3/4-3\epsilon-n\epsilon^2)F^2\ge F^2/2.
\]
When $\eta=1$, we have $e\xi/\chi\ge-1$.  When $\eta\ne1$ and
$e\xi<0$, we have $s\le|\dd\eta|_h$ and, by
\eqref{stretch:slope-bounds} and \eqref{stretch:eta-bounds},
\[
 \frac{|e\xi|}{\chi}
 \le\frac{|\dd\eta|_h s}{s_0\eta}
 \le\frac{|\dd\eta|_h^2}{s_0\eta}\le C_*.
\]
Using the last bound in Proposition~\ref{pressure:law}, we conclude that
\begin{align*}
 E&\ge\chi^2\bigl[(1/2-2C_*\epsilon)F^2-2C_*b_0V(t)\bigr]\\
  &\ge-C_2\chi^2V(t),\qquad C_2=2C_*b_0.
\end{align*}
The small metric modification now pays for this error.  Namely,
$\theta\chi^2\le2a^2s_0^2$ and
$b_*=-\beta\le\theta\le a^2s_0^2$, so
\[
 (V(t_u)-V(t))_+\le3a^2V(t),
                 \qquad V(t)\le2V(t_u).
\]
Fix $a\in(0,1/2]$ small enough that
\begin{equation}\label{stretch:a-choice}
 (2C_2+3)a^2\le d/4.
\end{equation}
Then the scalar error and the increase in reserve together cost at most
$(d/2)V(t_u)$.  This choice is made before taking $\rho$ small in the
locality argument.

Combining the two slope ranges gives
\begin{equation}\label{stretch:raw-estimate}
 \begin{split}
 \Dcal(h_u,t_u)\ge{}&\Dcal(h,t)+(V(t_u)-V(t))_+\\
 &-\delta(\tr_hh_0-\tr_{h_u}h_0)V(t_u)
                              -(d/2)V(t_u).
 \end{split}
\end{equation}
The inequality extends to boundary points of $\{s>0\}$ by smoothness.
On the interior of $\{s=0\}$ the metric and drift are unchanged, so it
holds there as well.  We have thus retained the full positive increase
of the reserve, apart from the trace loss and half the prescribed small
error.  The remaining half will absorb the final cutoff.

\subsection{Compact support and distance gain}
\label{stretch:finish}

Define the actual output on the original manifold by
\begin{equation}\label{stretch:output}
 h'=h+\zeta(h_u-h),\qquad t'=t+\zeta(t_u-t),
\end{equation}
extending the increments by zero.  The resulting data are smooth and
satisfy $h'\ge h$ and $t'\ge t$; their changes from $h,t$ have compact
support in $\{r_A<W\}$.
On the compact annulus
$\{.81W\le r_A\le.99W\}$, which contains every point where
$0<\zeta<1$, \eqref{stretch:locality} gives $C^3$ convergence
$u_\rho\to0$.  Consequently both the uncut and cut metric and drift
increments tend to zero in $C^2$ there.

To verify the inequality after cutoff without assuming any monotonicity
of $V$, consider its residual
\[
 \begin{split}
 \mathcal R(\widehat h,\widehat t)={}&
 \Dcal(\widehat h,\widehat t)-\Dcal(h,t)
                         -(V(\widehat t)-V(t))_+\\
 &+\delta(\tr_hh_0-\tr_{\widehat h}h_0)V(\widehat t)
                                      +dV(\widehat t).
 \end{split}
\]
By \eqref{stretch:raw-estimate},
$\mathcal R(h_u,t_u)\ge(d/2)V(t_u)\ge dT^2/2$.
The residual is continuous in the indicated $C^2$ data, including the
positive-part term.  Thus
$\mathcal R(h',t')-\mathcal R(h_u,t_u)\to0$ uniformly on the annulus.
For sufficiently small $\rho$, \eqref{stretch:conclusion} follows
there.  Outside this annulus, $\zeta$ is locally constant, with value
one or zero.  The data then agree locally with $h_u,t_u$, where
\eqref{stretch:raw-estimate} applies, or with the original $h,t$.
Completeness follows from $h'\ge h_0$.

It remains to quantify the distance gained.  Use the effective height
$\zeta u_\rho$, extended by zero.  By \eqref{stretch:height} and
\eqref{stretch:cutoffs}, it exceeds $W/5-1$ on $A$ and vanishes on $B$.
Decrease $\rho$ further so that
$|\dd(\zeta u_\rho)|_h\le.01$ throughout the cutoff transition.
Suppose a piecewise smooth curve from $A$ to $B$ had $h'$-length at most
$L_*W$.  Its $h$-length is also at most $L_*W$, and is at least $W$.
The absolute height variation in the cutoff transition, together with
that on $\{\zeta=1,s<.05\}$, is at most $.06L_*W$; repeated visits
are included by integrating the absolute derivative along the curve.
Where $\zeta=0$ there is no height variation.  Thus the remaining
portion, where $\zeta=1$ and $s\ge.05$, carries variation greater than
\[
 (1/5-.06L_*)W-1>.11W.
\]
On that portion $\mu=1$ and $h'=h+\dd u_\rho^2$.  Everywhere else
$h'\ge h$.  Applying the integral triangle inequality to the pair
consisting of the $h$-speed and the absolute height derivative restricted
to this portion gives
\[
 \operatorname{length}_{h'}(\gamma)
 >\sqrt{W^2+(.11W)^2}>L_*W,
\]
a contradiction.  This proves \eqref{stretch:distance-conclusion} and
completes the proof of Lemma~\ref{stretch:local}.
